\temaMundo
\chapter*{Do pós-guerra à lua}
\addcontentsline{toc}{chapter}{Do pós-guerra à lua}

\section*{Um mundo em reconstrução}

Quando os anos 1950 começaram, a Segunda Guerra Mundial havia terminado,
mas suas consequências ainda estavam presentes. Cidades destruídas,
fronteiras redesenhadas e milhões de vidas interrompidas faziam parte
da paisagem de um mundo que precisava reconstruir não apenas seus
territórios, mas também a esperança de que uma nova guerra pudesse ser
evitada.

Em 1945, enquanto o conflito chegava ao fim, representantes de dezenas
de países reuniram-se em São Francisco para construir uma nova
organização internacional. A Carta das Nações Unidas foi assinada em
26 de junho daquele ano e entrou em vigor em 24 de outubro, dando
origem oficialmente à Organização das Nações Unidas (ONU).%
\footcite{onu1945}

A ideia era ambiciosa: criar um espaço em que as nações pudessem
negociar, buscar soluções comuns e trabalhar pela manutenção da paz.
A experiência da guerra ainda era recente demais para ser esquecida.
O mundo que emergia daquele conflito, entretanto, logo descobriria
que a paz não significaria o fim das disputas.

\section*{Duas potências, dois mundos}

A reconstrução do pós-guerra aconteceu sob uma nova divisão de poder.
De um lado estavam os Estados Unidos e seus aliados; do outro, a
União Soviética e os países que orbitavam sua esfera de influência.

A Guerra Fria não seria uma guerra convencional entre as duas
superpotências. A disputa aconteceria também na política, na economia,
na ciência, na propaganda e na tecnologia. O planeta passou a viver
sob a possibilidade de um confronto que poderia ultrapassar as
fronteiras dos países diretamente envolvidos.

Enquanto adultos discutiam alianças, ideologias e armas, uma nova
geração crescia sob a promessa de que a ciência poderia transformar
o futuro. Aviões tornavam-se mais rápidos, a televisão começava a
entrar nas casas e os computadores, ainda enormes e raros, apontavam
para uma revolução tecnológica que poucas pessoas conseguiam imaginar.

Foi nesse ambiente de competição que surgiu uma nova fronteira:
o espaço.

\section*{A corrida espacial}

Em 4 de outubro de 1957, a União Soviética colocou em órbita o
Sputnik 1, o primeiro satélite artificial da Terra. Pequeno diante
da imensidão do espaço, o objeto produziu um impacto enorme na
imaginação mundial. Seus sinais de rádio podiam ser acompanhados
da Terra, transformando uma conquista tecnológica soviética em um
acontecimento acompanhado por pessoas muito além dos círculos
científicos.%
\footcite{nasa1957sputnik}

\begin{figure}[htbp]
    \centering
    \includegraphics[width=0.88\textwidth]{mundo/sputnik-1.jpg}
    \caption{Réplica do Sputnik 1 exibida no Cosmonautics Space Museum,
    Moscou.}
    \label{fig:sputnik1}

    \fonteImagem{GEAstro/UTFPR, ``Sputnik 1, o primeiro da era espacial'',
    2016.}
\end{figure}

A fotografia de uma réplica do Sputnik 1 ajuda a aproximar aquele
acontecimento distante de uma realidade concreta. O satélite que
parecia apenas um pequeno objeto metálico tornou-se um dos símbolos
da entrada da humanidade na era espacial.

Quatro anos depois, em 12 de abril de 1961, o soviético Yuri Gagarin
tornou-se o primeiro ser humano a orbitar a Terra, a bordo da
Vostok 1. Seu voo durou 108 minutos.%
\footcite{nasa1961gagarin}

A resposta norte-americana ganhou força especialmente depois da
decisão do presidente John F. Kennedy de estabelecer como objetivo
nacional levar um homem à Lua antes do final da década. A corrida
espacial deixava de ser apenas uma disputa por satélites e passava
a representar uma disputa pela capacidade de imaginar e realizar
o futuro.%
\footcite{nasa1969apollo11}

\section*{O caminho até a Lua}

O caminho até a Lua não aconteceu de uma só vez. Foi construído por
uma sucessão de experiências, fracassos e conquistas.

Em dezembro de 1968, a Apollo 8 levou seres humanos pela primeira
vez até a órbita lunar. No ano seguinte, a Apollo 11 completou a
etapa mais ousada desse processo. Lançada em 16 de julho de 1969,
a missão levou Neil Armstrong, Michael Collins e Edwin ``Buzz''
Aldrin ao espaço.

Em 20 de julho, Armstrong e Aldrin pousaram na Lua. Armstrong tornou-se
o primeiro ser humano a pisar em sua superfície, enquanto Collins
permaneceu em órbita lunar no módulo de comando.%
\footcite{nasa1969apollo11}

Centenas de milhões de pessoas acompanharam o acontecimento pela
televisão. Pela primeira vez, uma experiência vivida a centenas de
milhares de quilômetros de distância podia ser compartilhada
simultaneamente por uma parte enorme da humanidade.

Para uma criança que crescia naquele período, o espaço deixava de
ser apenas uma ideia distante. Ele passava a fazer parte das
conversas, das imagens, dos jornais, da televisão e dos sonhos.

\section*{Enquanto o mundo olhava para o céu}

Mas os anos 1950 e 1960 não foram apenas anos de descobertas e
conquistas tecnológicas.

Enquanto foguetes atravessavam a atmosfera, conflitos continuavam
a redesenhar a vida de milhões de pessoas. A Guerra Fria projetava
sua sombra sobre diferentes regiões do planeta, e a possibilidade
de um confronto entre as grandes potências permanecia presente.

Essa contradição marcou aquela geração. O mesmo período que levou
o primeiro ser humano ao espaço também foi atravessado por guerras,
disputas ideológicas e pelo medo de uma destruição em escala
desconhecida.

Era um mundo que aprendia a olhar para cima, para a Lua e para as
estrelas, enquanto ainda tentava compreender como viver em paz
sobre a própria Terra.

E foi nesse mundo que Maria começou sua vida.

Enquanto os grandes acontecimentos ocupavam jornais e telas de
televisão, a história de uma menina começava em um lugar muito
mais próximo e concreto: uma família, uma casa, uma pequena
localidade do Paraná e um cotidiano que, à primeira vista, parecia
muito distante da corrida espacial.

Mas o mundo não estava tão distante assim.

As mudanças econômicas, sociais e culturais daquele período também
chegariam ao interior do Brasil. E, dentro dele, chegariam à vida
de Maria.